\section{Coverage from shared continuation data}\label{sec:transfer}
The geometric argument provides joint residue entropy. Before choosing a
schedule of geometric operations, we prove the coverage principle that
will determine what the schedule must supply. The principle uses only
finite abelian residue groups and applies beyond Gaussian arithmetic.

\subsection{One shared conditioning cost}\label{sec:entropy}
We first isolate why a common data vector is useful. If one separately
conditioned on data chosen for each coordinate, its entropy cost would
be paid once per coordinate. A single vector can instead be paid for
once in the joint entropy, then distributed over the coordinates.

Let $\Lambda$ be a lattice, that is, a finitely generated free abelian
group, and let $k\ge1$ be an integer. For $1\le j\le k$ and
$\sigma\in\{+,-\}$ let
\[
 \phi_{j,\sigma}:\Lambda\longrightarrow G_{j,\sigma}
\]
be homomorphisms to finite abelian groups of order $p_j$.
Assume $T\le p_j\le2T$ and $T\ge2$, and put
\[
 \overline L=\frac1k\sum_{j=1}^k\log p_j,\qquad
 \av_c u_c=\frac1{2k}\sum_{j=1}^k\sum_{\sigma\in\{+,-\}}u_{j,\sigma}.
\]
Here $c=(j,\sigma)$ denotes one signed coordinate; write
$p_c=p_j$, $G_c=G_{j,\sigma}$, and $\phi_c=\phi_{j,\sigma}$.
For an integer $1\le s\le k$, choose $s$ distinct indices uniformly and one
independent fair sign at each. Denote this signed subset by $S$, and
write $F_S(Y)=(\phi_c(Y))_{c\in S}$ for a finite $\Lambda$-valued
variable $Y$. The choice of $S$ is auxiliary and independent of every
variable whose entropy is measured. In particular,
$\E_S H(F_S(Y))$ means the average of the entropies for fixed $S$.

\begin{lemma}[One shared conditioning cost]\label{lem:conditioning}
Let $(Y,\calH)$ be any pair of finite variables, with $Y$ taking values
in $\Lambda$, and let $S$ be selected independently as above. If
\[
 \E_S H(F_S(Y))\ge(1-\epsilon)s\overline L,
 \qquad H(\calH)\le\kappa s\overline L,
\]
then
\begin{equation}\label{eq:deficit-upper}
 \av_c[\log p_c-H(\phi_c(Y)\mid\calH)]
 \le(\epsilon+\kappa)\overline L.
\end{equation}
No independence between $Y$ and $\calH$ is required.
\end{lemma}
\begin{proof}
For each fixed $S$, the chain rule and subadditivity give
\[
 H(F_S(Y))-H(\calH)
 \le H(F_S(Y)\mid\calH)
 \le\sum_{c\in S}H(\phi_c(Y)\mid\calH).
\]
Each signed coordinate belongs to $S$ with probability $s/(2k)$.
Averaging and dividing by $s$ shows that the mean conditional entropy
is at least $(1-\epsilon-\kappa)\overline L$. The mean alphabet entropy
is $\overline L$, which proves the assertion.
\end{proof}

We will contradict \eqref{eq:deficit-upper} with the short-list bound
of Lemma~\ref{lem:list}. Its success event may depend on both the residue
and the data; independence of these variables is unnecessary.

\subsection{A transfer from later coverage}
We now mix a family of displacement laws. Each conditional law will
have a probability floor outside a small exceptional set. High joint
entropy of the starting position, together with a small cost for shared
data, will strengthen the exceptional-set bound for the mixture.

Let $A$ be a finite variable, let $y$ map its alphabet to $\Lambda$,
and put $Y=y(A)$. For every value $a$ of $A$, specify a finitely
supported displacement law $\nu_a$ on $\Lambda$. Conditional on $A$,
sample $h_1,\ldots,h_n$ independently with law $\nu_A$, where $n\ge1$
is an integer, and define
\[
 \calH=(h_1,\ldots,h_n),\qquad Z=Y+h_1.
\]
In the walk application, $A$ will record an exact intermediate time,
$Y$ its position, and $\nu_a$ the displacement law from that time to
the terminal position. Here $A$ is simply the finite mixing variable;
its influence on both $Y$ and $\calH$ is retained. The coordinate
subset $S$ used to measure entropy is independent of this whole experiment.

For a variable $X$ in an alphabet of size $p$, say explicitly that it
has $(\tau,\beta)$ coverage when
\begin{equation}\label{eq:coverage-definition}
 \#\{u:\Prb(X=u)<\tau/p\}<p^{1-\beta}.
\end{equation}
The strict inequality concerns the number of exceptional residues;
no claim is made that this set is empty.

The transfer below lowers the probability floor from $\tau'$ to $\tau$,
but improves the exceptional-set bound from $p^{1-\beta'}$ to
$p^{1-\beta}$, with $\beta\ge\beta'$. The exceptional sets of the
conditional laws may vary with $a$, so mixing them does not directly
give this improved bound. The theorem asks for three inputs: coverage
for each conditional law, a joint entropy budget that
can absorb the common data, and enough independent trials to make the
candidate lists small.

\begin{theorem}[One-step coverage transfer]\label{thm:transfer}
Use the signed coordinates and the experiment just defined, with
integers $1\le s\le k$ and $n\ge1$. Let
\[
 0<\tau\le3/8,\quad 0<\delta\le1/64,\quad
 \tau'=\tau+4\delta,\quad 0<\beta'\le\beta<1,
 \quad 0<\eta\le1,\quad 0\le\eta'\le1.
\]
Suppose that at every $a$ in the support of $A$, all but a fraction at
most $\eta'$ of the signed coordinates have $(\tau',\beta')$ coverage
for $\phi_c(y(a)+h_1)$ under the law conditional on $A=a$.
Assume also
\begin{align}
 \E_S H(F_S(Y))&\ge(1-\epsilon)s\overline L,
 &H(\calH)&\le\kappa s\overline L,\label{eq:transfer-entropy}\\
 \eta'&\le\eta\delta/100,
 &\epsilon+\kappa&\le\eta\delta\beta'/16,\label{eq:transfer-budget}
\end{align}
where $\epsilon,\kappa\ge0$. Finally, for every coordinate size $p=p_j$,
assume
\begin{align}
 p\exp(-\delta^2np^{-\beta}/2)&\le\delta/100,
 &\delta^{-1}&<p^{\beta'/2},\label{eq:transfer-size}\\
 \log(2e/\delta)&\le(\beta'/4)\log T.\label{eq:transfer-log}
\end{align}
Then all but a fraction at most $\eta$ of the coordinates have
$(\tau,\beta)$ coverage for $\phi_c(Z)$.
\end{theorem}

\begin{proof}
Call a coordinate deficient if its terminal law $\phi_c(Z)$ fails
\eqref{eq:coverage-definition}. Suppose their fraction is
$\theta>\eta$. For each deficient coordinate of size $p$ choose
\[
 K_c\subseteq\{u:\Prb(\phi_c(Z)=u)<\tau/p\},
 \qquad |K_c|=u_p:=\lceil p^{1-\beta}\rceil.
\]
Such a set exists, and $1\le u_p\le p$.
Using the same data $\calH$ for all coordinates, define
\begin{equation}\label{eq:list-definition}
 \calL_c(\calH)=\left\{x\in G_c:
 \sum_{\ell=1}^n\mathbf1_{\{x+\phi_c(h_\ell)\in K_c\}}
 \le n(\tau+\delta)u_p/p\right\}.
\end{equation}
We will show that the true residue belongs to this list often, but
that the list is small on average over deficient coordinates.

\paragraph{The true residue is often on the list.}
Every $Y+h_\ell$ has the same law as $Z$. The expectation of the count
in \eqref{eq:list-definition} at $x=\phi_c(Y)$ is therefore at most
$n\tau u_p/p$. Markov's inequality gives
\begin{equation}\label{eq:true-list}
 \Prb\{\phi_c(Y)\in\calL_c(\calH)\}
 \ge1-\frac{\tau}{\tau+\delta}
 =\frac\delta{\tau+\delta}\ge\delta.
\end{equation}
This step uses only the common one-trial marginal, not independence
of the trials before conditioning on $A$.

\paragraph{Coverage at the next checkpoint makes the list small.}
Fix $A=a$ and a coordinate whose conditional terminal law has
$(\tau',\beta')$ coverage. Translation by $\phi_c(y(a))$ shows that
the conditional law of $\phi_c(h_\ell)$ gives probability at least
$\tau'/p$ outside a set $E_{a,c}$ of size less than $p^{1-\beta'}$.
In any finite abelian group,
\[
 \sum_{x\in G_c}|(K_c-x)\cap E_{a,c}|=u_p|E_{a,c}|.
\]
Hence, except for at most $|E_{a,c}|/\delta$ candidates $x$,
\[
 |(K_c-x)\cap E_{a,c}|\le\delta u_p.
\]
For every other $x$, each conditional trial has hit probability at least
\[
 \tau'(1-\delta)u_p/p
 \ge(\tau+2\delta)u_p/p=:\lambda.
\]
Indeed $(\tau+4\delta)(1-\delta)-(\tau+2\delta)
=\delta(2-\tau-4\delta)>0$.

The count for such an $x$ stochastically dominates
$\operatorname{Bin}(n,\lambda)$, because its trials are independent
conditional on this exact $a$. The multiplicative lower-tail estimate
\cite[Exercise 8, p.~13]{Lugosi}, a standard concentration bound
(see also~\cite{BLM}), gives
\[
 \Prb\{\operatorname{Bin}(n,\lambda)
       \le n(\lambda-\delta u_p/p)\}
 \le\exp\!\left(-\frac{n(\delta u_p/p)^2}{2\lambda}\right)
 \le\exp(-\delta^2np^{-\beta}/2).
\]
For completeness, if $B\sim\operatorname{Bin}(n,\lambda)$ and
$0<v<1$, then
$\E e^{-tB}\le\exp(n\lambda(e^{-t}-1))$.
Markov's inequality with $t=-\log(1-v)$ bounds
$\Prb\{B\le(1-v)n\lambda\}$ by
$\exp[-n\lambda(v+(1-v)\log(1-v))]\le\exp(-n\lambda v^2/2)$.
Apply this with $v=\delta/(\tau+2\delta)$; the last inequality in the
displayed bound uses $u_p/p\ge p^{-\beta}$ and
$\tau+2\delta<1$.

A union bound over the at most $p$ candidates now shows that, with
conditional probability at least $1-\delta/100$, the list contains
only the exceptional candidates. By \eqref{eq:transfer-size}, their
number is less than
\[
 \delta^{-1}p^{1-\beta'}<p^{1-\beta'/2}.
\]
In particular,
\begin{equation}\label{eq:small-list}
 \Prb\{|\calL_c(\calH)|>p^{1-\beta'/2}\mid A=a\}
 \le\delta/100
\end{equation}
whenever this coordinate has the required coverage at $a$.

We have used next-checkpoint coverage only after fixing its exact value.
It remains to average over that value without assuming that a coordinate
deficient now is necessarily covered at the next checkpoint.

\paragraph{Averaging yields too much conditional information.}
At each $a$, the fraction of coordinates without the required coverage
is at most $\eta'$. Averaging first over $A$ and then over the
$\theta$ fraction of deficient coordinates, the mean probability of
such a failure is at most
\[
 \eta'/\theta<\delta/100.
\]
Let
\[
 a_c=\Prb\{\phi_c(Y)\in\calL_c(\calH),\quad
                  |\calL_c(\calH)|\le p_c^{1-\beta'/2}\}.
\]
Equations~\eqref{eq:true-list} and \eqref{eq:small-list} imply that its
mean $\overline a$ over the deficient coordinates satisfies
\[
 \overline a\ge\delta-\delta/100-\eta'/\theta\ge\delta/2.
\]
Lemma~\ref{lem:list} gives, for each of these coordinates,
\[
 \log p_c-H(\phi_c(Y)\mid\calH)
 \ge a_c(\beta'/2)\log p_c-h(a_c).
\]
Concavity of $h$ and $h(u)\le u\log(e/u)$ give
\[
 \av_{c\ \mathrm{deficient}}h(a_c)
 \le h(\overline a)
 \le\overline a\log(2e/\delta).
\]
After averaging and applying \eqref{eq:transfer-log}, we obtain
\[
 \av_{c\ \mathrm{deficient}}
 [\log p_c-H(\phi_c(Y)\mid\calH)]
 \ge\frac{\overline a\beta'}4\log T
 \ge\frac{\delta\beta'}8\log T.
\]
Every other conditional entropy deficit is nonnegative. Since
$\overline L\le\log(2T)\le2\log T$ and $\theta>\eta$, the mean over
all coordinates is strictly greater than
\[
 \frac{\eta\delta\beta'}8\log T
 \ge\frac{\eta\delta\beta'}{16}\overline L.
\]
This contradicts Lemma~\ref{lem:conditioning} and
\eqref{eq:transfer-budget}. Thus $\theta\le\eta$.
\end{proof}


The theorem has charged the repeated displacements once, through
$H(\calH)$. Its independence assumption is only conditional on the exact
checkpoint $A$. In the next section the schedule will make the joint
entropy error and this single data cost small enough to transfer coverage
at every step of a backward induction.
